\section{Substantive Comparability in Robust RAG}
\label{sec:robust-rag-landscape}

We use the literature map as a source of boundary cases, not as prevalence
evidence. Five \emph{operating questions} organise discovery: how evidence is
constructed, how a model uses it, how deliberate corruption changes the game,
how signals become actions, and how those actions are evaluated. They are a
different cut of the same literature from the four failure surfaces of
\S\ref{sec:taxonomy}, and both are different again from the five
\emph{reporting coordinates} of \S\ref{sec:coordinates}, which describe not a
literature but a single proposed contrast. The boundary cases below are what
made the coordinates necessary.

\paragraph{Evidence availability is not evidence use.}
Irrelevant-context studies test whether retrieved distractors redirect a
generator \citep{irrelevantContext2023}; noise-aware training changes the
generator's training distribution \citep{raat2024}; RGB supplies diagnostic
stress conditions \citep{rgb2023}; and Chain-of-Note changes how retrieved
documents are read \citep{chainOfNote2023}. A filter, a trained reader, and a
benchmark may share a harness, but their native scores do not rank
interchangeable interventions. Insufficient evidence creates another
boundary: when no passage can support an answer, re-retrieval or abstention
can dominate better distractor resistance \citep{sufficientContext2024}.

\paragraph{Support, authority, conflict, and attack are distinct.}
RAGTruth annotates unsupported spans, while evaluators estimate claim or
response support \citep{ragTruth2024,ragas2024,ares2024}; knowledge-conflict
methods instead study competition between parametric memory and retrieved
evidence \citep{tugOfWar2024,faithfulRAG2025,astuteRAG2025}. A response can be
correct but unsupported, or faithful to a weak source. Provenance and timestamp
determine which source should govern \citep{kilt2021,freshllms2024,
realtimeQA2023}. Deliberate corruption changes the game again: corpus
poisoning, indirect prompt injection, and backdoors differ in corpus write
access, query knowledge, retriever visibility, insertion budget, persistence,
and objective \citep{poisonedRAG2024,aip2025,phantom2024}, and their defenses
prevent, detect, certify, trace, quarantine, or recover under different
assumptions \citep{certifiedRAG2024,rrag017,rrag128}. An attack-success
comparison is transportable only within a matched capability envelope.

\paragraph{A signal is not a policy.}
Calibration predicts a property of a score \citep{guo2017calibration};
selective prediction trades coverage against risk
\citep{geifman2017selective}; sufficient-context classification
asks whether available passages can support an answer; adaptive retrieval
chooses whether to search, verify, answer, or abstain
\citep{calibRAG2025,sufficientContext2024,whenToRetrieve2024}. Self-RAG, CRAG,
and CtrlA expose this policy layer explicitly
\citep{selfRAG2023,crag2024,ctrla2025}. Comparing them requires a common
action set, budget, and loss: accuracy alone cannot price extra retrieval,
coverage, or refusal.

\paragraph{Comparison ladder.}
The strongest contrast reruns alternatives with the same corpus snapshot,
retriever, generator, evidence state, capability envelope, action set, and
outcome. When that is impossible, synthesis should retain only the common
scope supported by both studies. If a decision-critical field is absent, the
comparison is under-specified rather than evidence that either method is
better. Appendix~\ref{sec:extended-survey} gives paper-level contrasts and the
documentary-map boundary; the next section turns these recurring boundaries
into a claim-level audit.


\section{Extended Synthesis by Operating Question}
\label{sec:extended-survey}

Sections~\ref{sec:evidence}--\ref{sec:system} organise robust RAG by failure
surface, which is where a paper sits. This appendix cuts the same literature the
other way, by the \emph{operating question} a study answers, and develops the
paper-level contrasts the main text compresses. It is a synthesis, not an
exhaustive bibliography: the released \MapRows{}-row evidence map holds the long-tail
records, source strata and intervention labels.

Two limits carry over from \S\ref{sec:protocol} and govern how the contrasts
below may be read. The screening flow's first stage is documented but not
re-executable, so this synthesis identifies recurring scientific distinctions
rather than estimating prevalence or ranking evidence families; and the map
carries no study-level risk-of-bias weighting. A review-directed recall check
also surfaced older, pre-freeze publications absent from the frozen map, which
we cite as qualitative context and log as post-freeze discoveries without
entering them into any denominator. Where this survey states a number rather
than a distinction, it does so only for the frame of \S\ref{sec:applied}.

\subsection{When Does Retrieved Context Help?}

Irrelevant-context and noise-aware-training studies intervene at different
points. \citet{irrelevantContext2023} shows that retrieved distractors can
actively redirect generation and tests filtering and targeted fine-tuning.
\citet{raat2024} instead trains a generator across random, relevant, and
counterfactual retrieval conditions. RGB supplies stress conditions and
diagnostic abilities rather than a competing mitigation \citep{rgb2023}.
These papers can share an end-to-end harness, but their native results do not
rank a filter, a trained reader, and a benchmark as interchangeable objects.

Context construction creates a second boundary. Chain-of-Note changes how the
generator reads retrieved documents, whereas compression and reranking alter
what reaches it \citep{chainOfNote2023,rrag044,rrag058}. Long-context studies
show that retaining more passages can improve evidence availability while
worsening position sensitivity or distraction \citep{rrag039,rrag143}. A fair
comparison therefore fixes the candidate evidence, generator, decoding, and
budget when the intended claim is about the intervention rather than the
entire system.

Insufficient evidence is not merely severe distraction. Sufficient Context
labels whether the available passages can support an answer and evaluates
selective behavior \citep{sufficientContext2024}. A reader that resists a hard
negative and a controller that abstains when support is absent solve different
problems. Reports should state both answer-bearing support availability and
the scored action.

\subsection{Which Knowledge Should Govern an Answer?}

Knowledge-conflict work separates context--memory disagreement from conflict
among retrieved sources. Tug-of-War measures how models trade parametric and
contextual knowledge; FaithfulRAG models fact-level discrepancies; Astute RAG
reconciles imperfect retrieval with internal knowledge
\citep{tugOfWar2024,faithfulRAG2025,astuteRAG2025}. Although all three concern
``conflict,'' their experimental units, source assumptions, and desired
responses differ. The appropriate outcome may be answer correctness,
contextual faithfulness, calibrated disclosure, or source reconciliation.
Earlier entity-substitution experiments and PopQA isolate when parametric
memory resists or benefits from retrieved evidence
\citep{longpre2021conflicts,mallen2023trust}. ConflictBank broadens the
benchmark unit to millions of claim--evidence pairs across misinformation,
temporal conflict, and semantic ambiguity \citep{conflictBank2024}; its scale
does not make those conflict types interchangeable.

Grounding resources address another layer. RAGTruth provides span-level
hallucination labels, while RAGAS and ARES estimate response or context support
\citep{ragTruth2024,ragas2024,ares2024}. These evaluators can identify an
unsupported answer, but support alone does not establish that the cited source
is current or authoritative. A system can be faithful to weak evidence or
correct for a reason absent from its citations.

Source reliability and time therefore belong in the evidence record. Recent
work models source reliability directly \citep{sourceReliabilityRAG2025}, and
other studies examine source preference and contextual authority
\citep{rrag132,rrag147}. The comparison-relevant question is not only whether
the model used retrieved text, but whether the benchmark specified which
source should govern, at what timestamp, and under what corroboration rule.
KILT standardizes provenance against a fixed knowledge source, while FreshQA
and RealTimeQA test knowledge whose answer changes over time
\citep{kilt2021,freshllms2024,realtimeQA2023}. ALCE and generative-search
verifiability work separately score citation attribution and whether cited
sources actually support generated statements
\citep{alce2023,verifiability2023}.

Multi-agent methods add yet another intervention boundary. MAIN-RAG uses
agent voting to filter retrieved documents, MADAM-RAG assigns agents to
conflicting documents before aggregation, and DRAG introduces debate during
retrieval and generation \citep{mainRAG2025,madamRAG2025,drag2025}. General
multi-agent debate supplies methodological background
\citep{duMultiagentDebate2024}, but these RAG systems should be compared only
when the evidence allocation, agent budget, aggregation rule, and target
outcome match.

\subsection{What Changes Under an Attacker?}

Corpus poisoning makes evidence unreliability intentional. PoisonedRAG inserts
passages to steer target answers; earlier corpus-poisoning work manipulates
retrieval through adversarial passages
\citep{poisonedRAG2024,poisoningRetrievalCorpora2023}. Attack success is
conditional on corpus write access, insertion budget, query knowledge,
retriever visibility, and target scope. Omitting these fields can make two
attack rates numerically similar but scientifically incomparable.

Indirect prompt injection and backdoors target a different interface.
Retrieved text may contain instructions that override the user task, or a
trigger may activate behavior learned during training
\citep{aip2025,phantom2024,trojanRAG2024}. These attacks should not inherit
claims from factual poisoning merely because both place malicious text in the
context. Retrieval realism, instruction-following behavior, and trigger access
define distinct games.

Defenses likewise have different roles. Certified RAG gives guarantees under a
bounded corruption model; filtering and semantic isolation demote suspicious
documents; traceback identifies likely poisoned sources
\citep{certifiedRAG2024,rrag017,rrag128}. Prevention, certified accuracy,
detection, and recovery require different outcomes. Any cross-defense
comparison should pair the same capability envelope with clean utility and
state whether the defense prevents, detects, traces, quarantines, or recovers.

\subsection{How Does a Signal Become an Action?}

Calibration, answerability, and retrieval usefulness are related but not
identical signals. Calibration-oriented RAG predicts reliability, whereas
Sufficient Context asks whether the retrieved evidence can support an answer
\citep{calibRAG2025,sufficientContext2024}. A well-calibrated score can still
feed a poor threshold or action policy, and an insufficient-context detector
does not determine whether retrieval, clarification, or abstention has the
highest utility.

Adaptive systems make the policy explicit. Self-RAG learns when to retrieve
and critique; CRAG adds corrective retrieval; CtrlA separates uncertainty
sources to select control actions
\citep{selfRAG2023,crag2024,ctrla2025}. Their reported gains depend on the
available action set and resource budget. Answer-only accuracy cannot capture
the benefit or cost of extra retrieval, verification, or refusal.

Selective evaluation should therefore retain risk, coverage, action, and cost
rather than collapse them into one score. Trust-or-abstain methods, conformal
controllers, and uncertainty benchmarks instantiate different operating
points \citep{rrag002,rrag041,rrag001}. Comparisons are strongest when systems
share evidence and actions and are evaluated across a risk--coverage or
utility curve.

\subsection{What Can the Evidence Map Establish?}

The map documents that robust-RAG research spans evidence construction,
evidence use, adversarial integrity, and action control. Existing surveys
provide architectural, evaluation, trustworthiness, security, conflict, and
abstention views
\citep{ragSurvey2024,ragEvaluationSurvey2024,trustRAGCompass2024,
securityRAGSurvey2026,xuKnowledgeConflicts2024,wenAbstentionSurvey2025}.
Controlled searches for RAG best practices provide an empirical component
comparison under shared tasks \citep{ragBestPractices2024}, complementary to
survey taxonomies and to the claim-level audit studied here.
Table~\ref{tab:main-prior-work-comparison} states the narrower claim-level
increment tested here.

The map cannot establish that these families occur at representative rates,
that its taxonomy is complete, or that one family has stronger evidence.
Accepted-paper additions strengthen several historical lines, while
\MapPreprints{} of \MapRows{} mapped rows remain preprints against
\MapPeerReviewed{} peer-reviewed, and coding depth varies across legacy strata.
Map counts therefore provide documentary context only. Every measured claim in
this survey rests elsewhere: the taxonomy test on the constructed benchmark of
\S\ref{sec:tested}, the composition figures on the screened frame of
\S\ref{sec:applied}, and the design claim on the assignment-unit simulation.
